\begin{summarybox}
	\textbf{\textcolor{CSummary}{一句话核心}}
	对于焦点在 $x$ 轴的椭圆
	\[
		\frac{x^2}{a^2}+\frac{y^2}{b^2}=1\qquad (a>b>0),
	\]
	过焦点且斜率为 $k$ 的非垂直直线所截得的焦点弦长为
	\[
		|PQ|=\frac{2ab^2(1+k^2)}{a^2k^2+b^2}.
	\]

	\medskip

	\textbf{\textcolor{CSummary}{使用条件}}

	\begin{itemize}[leftmargin=*, itemsep=0.5em, topsep=0.4em]
		\item \textbf{条件1（焦点在 $x$ 轴）：}
		      椭圆方程必须为
		      \[
			      \frac{x^2}{a^2}+\frac{y^2}{b^2}=1\qquad (a>b>0),
		      \]
		      且
		      \[
			      c^2=a^2-b^2.
		      \]

		\item \textbf{条件2（直线过焦点）：}
		      直线必须经过一个焦点 $F(\pm c,0)$。
		      若直线不垂直于 $x$ 轴，可设斜率为 $k\in\mathbb R$ 并使用斜率公式；
		      若直线垂直于 $x$ 轴，则斜率不存在，应使用通径公式
		      \[
			      |PQ|=\frac{2b^2}{a}.
		      \]
	\end{itemize}

	\medskip

	\textbf{\textcolor{CSummary}{统一形式}}
	若过焦点的直线与 $x$ 轴夹角为 $\theta$，则
	\[
		|PQ|=\frac{2ab^2}{a^2\sin^2\theta+b^2\cos^2\theta}.
	\]
	其中：
	\[
		\theta=0 \quad\Longrightarrow\quad |PQ|=2a,
	\]
	对应长轴；
	\[
		\theta=\frac{\pi}{2} \quad\Longrightarrow\quad |PQ|=\frac{2b^2}{a},
	\]
	对应通径。

	\medskip

	
    \textbf{\textcolor{CSummary}{考前速记：两个特殊值 + 两条恢复路线}}
    \[
        \underbrace{\theta=0\ \Longrightarrow\ |PQ|=2a}_{\text{水平长轴，最长}},\qquad
        \underbrace{\theta=\frac\pi2\ \Longrightarrow\ |PQ|=\frac{2b^2}{a}}_{\text{竖直通径，最短}}.
    \]
    \textbf{更易记的斜率式：}
    \[
        \boxed{|PQ|=\frac{2a(1+k^2)}{1+\frac{a^2}{b^2}k^2}}
        \quad(k\in\mathbb R).
    \]
    \textbf{最快恢复（已知焦半径）：}
    \[
        PF_2=a-\frac ca x_P,\quad x_P=c+PF_2\cos\theta
        \ \Longrightarrow\ PF_2=\frac{b^2}{a+c\cos\theta}.
    \]
    对向半弦只需变号：$QF_2=\frac{b^2}{a-c\cos\theta}$，再相加。

    \textbf{从零恢复（忘了焦半径）：}
    $PF_1+PF_2=2a$，在焦点三角形内用\textbf{余弦定理}，
    得同一个 $\frac{b^2}{a+c\cos\theta}$；另一端换方向后相加。
    \textbf{核心动作：}一条焦点弦 $=$ 两条焦半径之和。

    \medskip


\textbf{\textcolor{CSummary}{关键提醒}}

	\begin{itemize}[leftmargin=*, itemsep=0.5em, topsep=0.4em]
		\item \textbf{易错点（焦点位置）：}
		      若椭圆焦点在 $y$ 轴，原斜率公式不能直接套用，需要使用对应变式。

		\item \textbf{检查项（直线条件）：}
		      务必确认直线是否经过焦点；不过焦点的弦不要使用本公式。

		\item \textbf{边界项（垂直直线）：}
		      直线垂直于 $x$ 轴时斜率不存在，不能硬代入斜率公式，应单独用通径公式。
	\end{itemize}
\end{summarybox}